\section{H2: Order imbalance and pressure}\label{sec:pressure}

\subsection{Order imbalance in the window}

Table~\ref{tab:pressure} reports the order imbalance of Section~\ref{sec:hypotheses} for
each type of session. Trades are classified as buyer- or seller-initiated by the side whose
order arrived later, which is the side that crossed the spread.

On ordinary sessions the imbalance over the window is small in absolute value,
\PressAbsoibLiquidNormal{} in the liquid group and \PressAbsoibPlaceboNormal{} in the placebo
group, and it has no systematic sign. On monthly expiries the absolute imbalance of the liquid
group rises by \PressAbsoibLiquidMonthlyexpiry{} (standard error
\PressAbsoibLiquidMonthlyexpirySe{}), and its average changes by $\PressOibLiquidMonthlyexpiry{}$, toward selling. Neither change is large relative to the variation across
sessions, and neither differs significantly from the placebo group. The window is not
dominated by one side of the market on expiry days.

The imbalance does, however, become more persistent. On ordinary sessions the imbalance of one
five-minute block carries over to the next with a coefficient of
\OibdynPersistenceLiquidLag{} in the liquid group. On monthly expiries the coefficient rises by
\OibdynPersistenceLiquidLagxmonthlyexpiry{} (standard error
\OibdynPersistenceLiquidLagxmonthlyexpirySe{}), so that buying or selling that begins early in
the window continues through it. The rise is visible in all three groups but is largest in
the liquid group; relative to the placebo group it is \OibdynPersistenceDidliquidLagxmonthlyexpiry{}
(standard error \OibdynPersistenceDidliquidLagxmonthlyexpirySe{}), which is not significant.
On weekly index expiries persistence does not change. Figure~\ref{fig:oib} shows the absolute
imbalance block by block and the persistence coefficients.

\begin{figure}[htbp]
  \centering
  \includegraphics[width=\textwidth]{oib.pdf}
  \caption{Left: mean absolute order imbalance in each five-minute block of the window, for the
    liquid and placebo groups on monthly expiries and other sessions. Right: persistence of the
    imbalance from one block to the next, with 95\% intervals.}
  \label{fig:oib}
\end{figure}
 Persistent imbalance is what the
execution of large orders through the window produces, whether they are placed to close
expiring positions or to move the price. The move from the window's first minute to its last
is larger on monthly expiries, by \PressAbsmoveOneFiveThreeZeroLiquidMonthlyexpiry{} basis
points in the liquid group, and a given imbalance moves the price by
\OibdynPriceDerivativeOib{} basis points per unit of imbalance on any session, with no
significant change on expiry days.

\begin{table}[htbp]
\centering
\small
\caption{Order imbalance in the settlement window. Order imbalance is
$(BI - SI)/(BI + SI)$, where $BI$ and $SI$ are buyer- and seller-initiated volume. The first
three panels report coefficients on indicators for the session type with security fixed
effects; the last column is the difference between the securities with derivatives and the
placebo group, with security and session fixed effects. The last panel regresses the imbalance
of each five-minute block on that of the previous block, with interactions for the session
types. Standard errors clustered by session in parentheses.}
\label{tab:pressure}
\resizebox{\textwidth}{!}{\input{generated/tables/pressure}}
\end{table}

\subsection{Pressure or congestion in the order book}

A window that is busier because more participants trade in both directions and a window in
which the price is pushed in one direction both show wider spreads and more cancellations.
They differ in three respects. Under sustained pressure, returns over successive seconds lean
the same way and the variance ratio rises; signed order flow is persistent; and a larger part
of the price impact of order flow is permanent. Under heavier two-sided trading the variance
ratio falls, flow is no more persistent, and more of the impact is transitory.

None of the three measures moves toward pressure. The variance ratio is
\VarianceRatioExpiry{} on expiry sessions against \VarianceRatioControl{} on controls, the
persistence of signed flow is unchanged, and the permanent share of price impact is
\ImpactPermanentShareExpiry{} against \ImpactPermanentShareControl{}. Against the placebo group
two of them move in the direction of congestion: relative to the placebo group, the change in the
variance ratio of securities with derivatives is $\DidVarianceratioPlacebo{}$
($p = \DidVarianceratioPlaceboP{}$), and the change in the permanent share of price impact
$\DidPermanentsharePlacebo{}$ ($p = \DidPermanentsharePlaceboP{}$). The price impact of a given flow over ten seconds rises,
from \ImpactShortBpsControl{} to \ImpactShortBpsExpiry{} basis points, but less of it lasts.
Figure~\ref{fig:pressure} shows the distributions.

\begin{figure}[htbp]
  \centering
  \includegraphics[width=\textwidth]{pressure_measures.pdf}
  \caption{The three measures that distinguish directional pressure from heavier two-sided
    trading, expiry against control sessions. Under sustained pressure the variance ratio and
    the persistence of flow rise and a larger share of impact is permanent.}
  \label{fig:pressure}
\end{figure}

\subsection{Whose flow moves the settlement price?}

Table~\ref{tab:flow} relates the window's move and the next morning's return to the net
aggressor flow in the window. Across the securities with derivatives, net aggressive buying of
one per cent of median session volume raises the settlement price by
\FlowDriftAllDerivativeFlow{} basis points on an ordinary session, and the effect on monthly
expiries is not different.

The participant categories differ sharply. Aggressive flow of custodians and of other clients
moves the settlement price in the direction of the flow, by \FlowDriftParticipantDerivativeFlowcustodian{}
and \FlowDriftParticipantDerivativeFlowncnp{} basis points per per cent of volume. Aggressive
flow of proprietary traders is associated with a move in the opposite direction,
$\FlowDriftParticipantDerivativeFlowproprietary{}$ basis points, and is followed by a
next-morning return of \FlowNextrevParticipantDerivativeFlowproprietary{} basis points in its
own favour. Proprietary traders thus trade aggressively against the window's move and are
rewarded overnight, which is the behaviour of liquidity supply by arbitrageurs rather than of
price pressure. None of these relations changes on monthly expiries: no participant category's
window flow has a larger price effect, or is followed by a larger reversal, on the days on
which the settlement is struck.

\begin{table}[htbp]
\centering
\small
\caption{Net aggressor flow in the window, the window's move and the next morning's return.
Flow is the volume bought by aggressive buyers less the volume sold by aggressive sellers, in
per cent of the security's median session volume, in total and by the aggressor's category.
Regressions include security and session fixed effects and the session-type indicators and
their interactions with flow; only the flow terms and their interaction with the monthly
expiry are shown. Standard errors clustered by session in parentheses.}
\label{tab:flow}
\resizebox{\textwidth}{!}{\input{generated/tables/flow}}
\end{table}
